\documentclass[11pt,a4paper]{article}
\usepackage[margin=26mm]{geometry}
\usepackage{amsmath,amssymb,amsthm,hyperref}
\newtheorem{theorem}{Theorem}\newtheorem{proposition}[theorem]{Proposition}
\newcommand{\gen}[1]{\langle #1\rangle}
\title{Three elementary interfaces for a survey of string C-groups}
\author{Human Compression project}\date{}
\begin{document}\maketitle
This note treats three precise statements in the fixed survey edition
arXiv:1910.08843v2 of Leemans~\cite{survey}. It extends the acquired elementary
proof of the simple-group intersection lemma to the full hypothesis stated
in survey Theorem~3.2; it derives the symmetric rank intervals with the correct
endpoint type; and it separates the valid projective form obstruction from
the stronger, false wording of survey Theorem~5.7. These are bounded components,
not a disposition of the survey's other classifications, tables or questions.

\section{The rank-three intersection argument}
\begin{theorem}[Conder--Oliveros]\label{three}
Suppose a nontrivial finite group $G$ has no nontrivial cyclic normal subgroup
and is generated by involutions $a,b,c$, where $a$ and $c$ commute. Then
$(a,b,c)$ is a rank-three string C-group marking.
\end{theorem}
\begin{proof}
A group generated by two involutions is dihedral and has a nontrivial cyclic
normal subgroup, including in the abelian cases of order two or four. Thus
our three generators are distinct and independent. Put
$A=\gen{a,b}$, $B=\gen{b,c}$ and $H=A\cap B$. Both $A$ and $B$ are proper
subgroups. Assume that $H>\gen b$. Since $b$ is a reflection in each of the
two dihedral groups, $H$ is dihedral. Its respective rotation subgroups are
$X=H\cap\gen{ab}$ and $Y=H\cap\gen{bc}$, both of index two in $H$.
Unless $H$ is a four-group its cyclic rotation subgroup is unique, so $X=Y$.
The same conclusion may also hold when $H$ is a four-group. In either case
$X=Y\ne1$ is normal in both $A$ and $B$: every subgroup of the cyclic
rotation subgroup of a dihedral group is normal. It is therefore a nontrivial
cyclic normal subgroup of $G$, a contradiction.

In the remaining case write $H=\{1,u,v,b\}$ with $X=\gen u$, $Y=\gen v$
and $b=uv$. Then $u$ is central in $A$ and $v$ is central in $B$. Every
one of $a,v$ commutes with every one of $c,u$, using also $[a,c]=1$ and
$[u,v]=1$. Consequently
\[
 G=UV,\qquad U=\gen{a,v},\quad V=\gen{c,u},\quad [U,V]=1.
\]
The dihedral group $U$ is normal in $G$, and its commutator subgroup is
cyclic and characteristic in $U$. Our hypothesis forces $U'=1$; similarly
$V'=1$. Hence $G$ is abelian, again contradicting the absence of a
nontrivial cyclic normal subgroup. We have proved $A\cap B=\gen b$.

The subgroup $C=\gen{a,c}$ is a four-group. Its intersection with $A$
is exactly $\gen a$: a larger intersection would contain $c$ and imply
$G=A$. Similarly $C\cap B=\gen c$. These three intersections and
independence give every subset intersection, including the empty subset.
\end{proof}
The source attributes this result to Conder and Oliveros, Corollary~4.2.
The proof reuses the actual dihedral argument in the accepted L011 application;
only its final simple-group contradiction is replaced by the characteristic
commutator argument above. Smooth quotients of the string Coxeter group
$[k,m]$ meet the hypotheses on their three distinguished generators.

\section{Symmetric groups: the endpoint and all intermediate ranks}
The acquired contraction theorem states that an irreducible string C-marking
$(r_0,\ldots,r_{R-1})$ of rank $R\ge4$ can be replaced by
$(r_1,r_0r_2,r_3,\ldots,r_{R-1})$ when
$r_0\in\gen{r_0r_2,r_3}$. If $r_2r_3$ has odd order $m$, this membership
follows from $((r_0r_2)r_3)^m=r_0$. The complete theorem and its intersection
proof are acquired from the accepted common contraction paper, not assumed
from the conclusion to be recovered here.

\begin{proposition}\label{sym}
For $n\ge4$, every rank $3\le r\le n-1$ occurs for $S_n$. At rank $n-1$
the simplex marking has type $\{3^{n-2}\}$. For $3\le r\le n-2$ a marking
has type $\{n-r+2,6,3^{r-3}\}$.
\end{proposition}
\begin{proof}
Start with the adjacent-transposition marking of $S_n$. A selected edge set
in its path generates the direct product of symmetric groups on the resulting
intervals. Intersecting two such subgroups gives the group of their common
interval refinement, exactly the subgroup on the common selected edges.
This proves IP. All adjacent products have order three.

After $s\ge1$ contractions, where $r=n-1-s\ge3$, the permutation graph is
still the path on $n$ vertices. Its first $s+2$ edges alternate in labels
$0,1$, ending in label $1$; its remaining edges have labels
$2,3,\ldots,r-1$. This follows inductively from the replacement: the two
old alternating labels are exchanged, and the old label~2 becomes the next
edge of the new label~1. The labels from~3 onward decrease by one.
At every step the bond tested by contraction joins two consecutive pure
path transpositions, so has order three. Generation and IP are therefore
preserved by the acquired theorem.

The product of the first two marks is a cycle on the $s+3$ initial vertices:
alternating along the path gives one orbit and visits every vertex once.
Its order is $s+3=n-r+2$. The second mark has at least two transpositions.
Its last transposition meets the pure label-2 edge; their product is a
3-cycle, while at least one disjoint transposition of that second mark
remains. The adjacent product therefore has order six. Every subsequent
adjacent product has order three. This proves the displayed type and the
proposition.
\end{proof}
This recovers survey Theorems~4.1 and~4.3, except that the type in the latter
must be restricted to $r\le n-2$. Its literal endpoint assertion is already
false at $n=4,r=3$: the proposed type is $\{3,6\}$, but $S_4$ has no element
of order six. Indeed its disjoint cycle lengths have sum four, so neither a
6-cycle nor a disjoint 3-cycle and 2-cycle is available. The rank-existence
conclusion is unaffected. Attribution remains with Fernandes--Leemans for
the original construction and Brooksbank--Leemans for contraction.

\section{What the three-dimensional obstruction actually asserts}
Survey Theorem~5.7 states that any $G\le\mathrm{GL}_3(q)$ admitting a string
C-group marking has odd characteristic and satisfies
$\Omega(V,f)\le G\le I(V,f)$ for a nondegenerate symmetric form~$f$.
Both the missing irreducibility condition and the lower containment matter.
They are logically separate defects.

Let $S_4$ permute the coordinates of $F^4$, and put
\[
 W=\{(x_1,x_2,x_3,x_4):x_1+x_2+x_3+x_4=0\}.
\]
Use the ordered adjacent transpositions $(12),(23),(34)$, restricted to $W$.
For $F=\mathbb F_2$ or $\mathbb F_5$ the action is faithful. To see this,
if a permutation fixes every $e_i-e_4$, it fixes each unordered support
$\{i,4\}$, so fixes their common point~4 and then each point~$i$.
Thus the restricted tuple is still the simplex string C-marking of $S_4$,
of rank three and type $\{3,3\}$. Over $\mathbb F_2$ this immediately
contradicts the unqualified odd-characteristic assertion.

Over $\mathbb F_5$ the action on $W$ is absolutely irreducible. Indeed work
over any extension field and take a nonzero vector $x$ in an invariant
subspace of $W$. Its coordinates cannot all agree, because $4\ne0$ in
characteristic five and their sum is zero. Choose $i,j$ with $x_i\ne x_j$.
Then $x-x^{(ij)}=(x_i-x_j)(e_i-e_j)$ belongs to the subspace. The $S_4$-orbit
of this vector spans $W$. Hence that subspace is all of $W$.
But $|S_4|=24$, whereas every nonsingular three-dimensional orthogonal
form over $\mathbb F_5$ has $|\Omega_3(5)|=5(5^2-1)/2=60$.
Consequently the claimed lower containment is impossible even after adding
absolute irreducibility. The standard orthogonal order formula is the only
classical group input in this last comparison.

The cited Brooksbank--Vicinsky author version of 28 June 2009 explicitly
imposes absolute irreducibility, but its introductory theorem also prints
this lower containment. Theorem~3.3 in its body states form preservation,
without this lower containment.
We do not treat the printed containment as a valid retained premise.
The exact theorem we acquire instead is the proved L011 projective form
obstruction: an irreducible string C-group in $\mathrm{PGL}_3(F)$ has rank
at most four, odd characteristic, and a linear lift preserving a nondegenerate
symmetric form. It does not assert containment of the full orthogonal
derived group.

For completeness this acquired result does imply the survey's stated
exclusion of both $\mathrm{PSL}_3(q)$ and $\mathrm{PGL}_3(q)$. Their natural
projective actions are irreducible. In even characteristic the form obstruction
excludes them. For odd $q\ge3$, a lifting group projecting onto either must
have order at least $|\mathrm{PSL}_3(q)|$, but is contained in $O_3(q)$ and
\[
 \frac{|O_3(q)|}{|\mathrm{PSL}_3(q)|}
 =\frac{2\gcd(3,q-1)}{q^2(q^3-1)}<1.
\]
This is impossible. The order formulas and irreducibility of the natural
classical action are the same retained inputs as in the accepted L011 proof.
No inference that the projective image of an arbitrary linear C-marking
again has IP is used.

The finite companion merely checks the two displayed $S_4$ actions and
small instances of the path formulas. The uniform arguments above, and the
complete acquired contraction and projective-form proofs, are the proof
premises. None of the survey's open questions is answered by this note.
\begin{thebibliography}{9}
\bibitem{survey} D. Leemans, \emph{String C-group representations of almost
simple groups: a survey}, arXiv:1910.08843v2, 11 May 2020, pp.~4,6,12.
\bibitem{CO} M. Conder and D. Oliveros, \emph{The intersection condition for
regular polytopes}, J. Combin. Theory A 120 (2013), 1291--1304.
\bibitem{BL} P. A. Brooksbank and D. Leemans, \emph{Rank reduction of string
C-group representations}, Proc. Amer. Math. Soc. 147 (2019), 5421--5426.
\bibitem{BV} P. A. Brooksbank and D. A. Vicinsky, \emph{Three-dimensional
classical groups acting on polytopes}, Discrete Comput. Geom. 44 (2010),
654--659; author version 28 June 2009.
\bibitem{acquired} Human Compression project, \emph{Contraction of string
markings through a central factor}; \emph{Projective involutions and support
spaces in small dimension}; \emph{The complete exclusion list for simple
string C-groups}. Exact accepted editions and reviews are bound in the
companion dependency record.
\end{thebibliography}\end{document}
